\begin{proposition}[B.4.16; Proved; scope of the insertion operations]
\label{push:B416}
The local operation B.4.15 attaches a triangulated
cylinder to the original twelve-cycle in the dual
complex and admits an elementary collapse back to that
cycle. In a compact normal-crossings model, performing
it at a regular place consequently preserves the
homotopy type of the global dual complex.

The same conclusion holds for finite sequences of these
insertions, twists by a connected holomorphic group
action preserving the prescribed normal-crossings
model, and toroidal subdivisions of its central cone
complex. In particular, starting with a model whose
dual complex has a nonzero first integral cohomology
class, these operations cannot produce a semistable
$S^6$ centre. They also cannot change the component
rotations in B.3.1.

This statement concerns the specified generated class
of operations. It makes no assertion that arbitrary
analytic caps or arbitrary changes of the punctured
comparison belong to that class.
\end{proposition}

\begin{proof}
Write $H_i$, $i\in\Z/12$, for the original component
vertices and $V_i$ for the new vertical ones. At a node
between an even component $i$ and its odd neighbour
$j$, the two charts (B.4.15.3) have triple strata
\[
 \{V_i,V_j,H_j\},\qquad
 \{V_i,H_i,H_j\}.
 \tag{B.4.16.1}
\]
The square with horizontal edges $H_iH_j$ and $V_iV_j$
is thus divided by the diagonal $V_iH_j$. Over all
twelve nodes, (B.4.16.1) is exactly a triangulation of
$C_{12}\times[0,1]$, attached along the old horizontal
cycle $C_{12}\times\{0\}$.

There is a literal elementary collapse relative to the
old cycle. First remove each triangle containing a top
edge $V_iV_j$, paired with that free edge. Next remove
each remaining triangle, paired with its now free
diagonal $V_iH_j$. Finally remove each vertical edge
$H_iV_i$, paired with its free vertex $V_i$. What remains
is precisely the original twelve-cycle. These are
$36$ elementary collapses. Equivalently, the retraction
sends $V_i$ to $H_i$ and is the identity on all old
vertices and edges.

This description is valid inside a larger dual complex
$K$. The only new strata lie over the base origin,
so the new complex is the mapping-cylinder attachment
\[
 K\ \cup_{\,f:C_{12}\to K}\
           (C_{12}\times[0,1]).
\]
It retracts onto $K$, regardless of possible global
identifications among the local horizontal branches.
For self-crossings one uses the ordered local-branch
dual complex, so no identification of distinct local
branches is silently discarded.

A connected group acts trivially on the finite sets
of component labels and connected components of every
stratum: its permutation is locally constant and is
the identity at the identity element. Cartier twists
by such an action therefore keep this same incidence
complex. Toroidal subdivisions subdivide the
corresponding central polyhedral complex and hence
preserve its underlying homotopy type. These arguments
apply successively, proving the assertion for finite
compositions. Around each regular insertion point the
relative component monodromy is still the identity;
no peripheral rotation is added to B.3.1.

For completeness, a compact normal-crossings space $D$
has a natural injection
\[
 H^1(\Delta(D),\C)\ \lhook\joinrel\longrightarrow\
 H^1(D,\C),
 \tag{B.4.16.2}
\]
where $\Delta(D)$ denotes its local-branch dual
complex with connected strata. Use the ordered finite
branch spaces $D^{[p]}$ and take the alternating
summand for the branch-permutation action. Over $\C$
this is an exact direct-summand operation. Resolve the
constant sheaf by this augmented alternating branch
complex. At a point its row is the augmented simplex
complex, so it is exact, including for self-crossings.
The associated first-quadrant spectral sequence has
\[
 E_1^{p,q}=H^q(D^{[p]},\C)^{\mathrm{alt}},
 \qquad
 E_2^{1,0}=H^1(\Delta(D),\C).
\]
No higher differential enters or leaves $(1,0)$.
Its filtration piece in total degree one gives
(B.4.16.2). Equivalently one can use the orientation
local system on each unordered branch stratum; an
ordinary constant sheaf there would be incorrect if
branches permute. A nonzero integral first cohomology
class of the finite dual complex remains nonzero over
$\C$, because first integral cohomology is torsion-free.
Thus a surviving circle class forces $H^1(D,\C)\ne0$.
By B.4.2 every semistable centre of
a simply connected nearby $S^6$ is simply connected,
which is incompatible with this conclusion.

The finite incidence and collapse certificates themselves
are integral. They are
saved in \texttt{code/ab\_push\_regular\_cap\_dual.py}
and its output. For the twelve-cycle the chain ranks
are $24,48,24$, the nonzero Smith entries of the two
boundary maps are respectively $1^{23}$ and $1^{24}$,
and the integral homology is $(\Z,\Z,0)$.
The explicit retraction pulls the integral angular
cocycle back with winding one. These exact checks
certify the finite combinatorics; the preceding chart
calculation supplies its analytic interpretation.
\end{proof}
